\documentclass[11pt,a4paper]{article}
\usepackage{notes}
\usepackage{fancyhdr}

\pagestyle{fancy}
\fancyhf{}
\lhead{Geophysics IIIA: Potential Fields \& Geothermics}
\rhead{Week 8 Geotherm Curvature}
\cfoot{\thepage}

\title{Week 8 --- Activity: Energy Conservation \\ and Geotherm Curvature}
\author{Geophysics 3A: Potential Fields \& Geothermics}
\date{\today}

\begin{document}
\maketitle

\section*{Purpose}
To understand how crustal heat production alters lithospheric temperatures even when surface heat flow is unchanged.

\section*{Scenario}
Two continental terranes have identical surface heat flow, $q_0$, but different
amounts of heat supplied from the mantle at the base of the crust.

\section*{Tasks}
\begin{enumerate}
  \item Write the steady-state energy balance for the crust:
  \[
    q_0 = q_b + \int_0^{z_c} A(z)\,dz
  \]

  \item If Terrane A has a larger basal heat flow ($q_b$) than Terrane B, which must
  have higher integrated crustal heat production? Explain.

  \item Consider the steady-state geotherm:
  \[
    T(z) = T_0 + \frac{q_0}{k}z - \frac{A}{2k}z^2
  \]
  Identify which term controls curvature.

  \item Without performing calculations, explain how greater curvature affects:
  \begin{itemize}
    \item temperature at the Moho,
    \item depth of intersection with the mantle adiabat.
  \end{itemize}
\end{enumerate}

\section*{Key Insight}
Explain why surface gradients alone are insufficient to determine deep crustal
temperatures.

\end{document}